% Compile twice with XeLaTeX: xelatex -interaction=nonstopmode -halt-on-error EXTENSIONS_REPORT.tex
\documentclass[10pt,a4paper]{scrartcl}

\usepackage[margin=20mm]{geometry}
\usepackage{fontspec}
\usepackage{mathtools}
\usepackage{unicode-math}
\usepackage{microtype}
\usepackage{amsthm}
\usepackage{booktabs}
\usepackage{tabularx}
\usepackage{array}
\usepackage{enumitem}
\usepackage{xcolor}
\usepackage{hyperref}
\usepackage{bookmark}
\usepackage{cleveref}
\usepackage{xurl}
\usepackage{scrlayer-scrpage}

\setkomafont{section}{\Large\bfseries}
\setkomafont{subsection}{\large\bfseries}
\setkomafont{disposition}{\normalfont\bfseries}
\setlength{\parindent}{0pt}
\setlength{\parskip}{0.55em}
\setlist[itemize]{nosep,leftmargin=1.5em}
\setlist[enumerate]{nosep,leftmargin=1.7em}
\renewcommand{\arraystretch}{1.16}
\emergencystretch=2em

\hypersetup{
  unicode=true,
  colorlinks=true,
  linkcolor=blue!45!black,
  urlcolor=blue!45!black,
  pdftitle={Ten Further Extensions of the Prime-GPF Theory},
  pdfauthor={Anders Hellstrom}
}

\pagestyle{scrheadings}
\clearpairofpagestyles
\ihead{Prime-GPF theory extensions}
\ohead{September 2026}
\cfoot{\pagemark}

\newtheorem*{remark}{Remark}
\newcommand{\PP}{\mathbb P}
\newcommand{\card}[1]{\#\!\left(#1\right)}
\newcommand{\ind}{\mathbf 1}
\newcommand{\BigO}{\mathcal O}

\title{Ten Further Extensions of the Prime-GPF Theory}
\subtitle{Mathematical report and Lean formalization status}
\author{Anders Hellstrom}
\date{15 September 2026}

\begin{document}
\maketitle
\tableofcontents
\bigskip

This report records ten mathematical extensions of the three prime-GPF operations
\[
A(p,q)=P^+(p+q+1),\qquad
M(p,q)=P^+(pq+1),\qquad
E_p(q)=P^+(p^q+1),
\]
with prime inputs and anchors. The mathematical proofs are given below. Formal Lean status is reported separately and is deliberately stricter than mathematical status.

\section*{Formalization status}
\addcontentsline{toc}{section}{Formalization status}

\begin{center}
\small
\begin{tabularx}{\textwidth}{@{}>{\centering\arraybackslash}p{0.10\textwidth}p{0.18\textwidth}X@{}}
\toprule
\textbf{Extension} & \textbf{Mathematical status} & \textbf{Lean status on \texttt{extensions-5-6-8-9-10-20260915}} \\
\midrule
1 & proved below & not yet formalized \\
2 & proved below & not yet formalized \\
3 & proved below & not yet formalized \\
4 & proved below & not yet formalized \\
5 & proved below & core escape and visit-count statements compiler-verified \\
6 & proved below & finite image bound (9) and deep-image classification (11) compiler-verified; asymptotic ratio (10) not yet formalized \\
7 & proved below & not yet formalized \\
8 & proved below & inequalities (14)--(16) compiler-verified; separate prime-polynomial corollary added, compiler validation pending \\
9 & proved below & Lean proof present; not yet compiler-verified (current branch CI fails in \texttt{Extension9.lean}) \\
10 & proved below & algebraic separation compiler-verified; dynamical separation not yet formalized \\
\bottomrule
\end{tabularx}
\end{center}

The verified Lean material is in \path{PrimeGPF/Extensions.lean}; extension 9 is developed in \path{PrimeGPF/Extension9.lean}.

\section{Counting lemma}

For a finite nonempty set of primes $T$, put
\[
d=|T|,\qquad C_T=\frac1{d!\prod_{\ell\in T}\log\ell},
\]
and let $S_T(Y)$ count positive integers at most $Y$ whose prime factors belong to $T$. Then
\begin{equation}
S_T(Y)=C_T(\log Y)^d+O_T((1+\log Y)^{d-1}). \tag{*}
\end{equation}

\begin{proof}
Unique factorization identifies the counted integers with lattice points
\[
e_\ell\ge0,\qquad \sum_{\ell\in T}e_\ell\log\ell\le\log Y.
\]
The simplex has volume $C_T(\log Y)^d$. Place a unit cube at each lattice point. Their union is sandwiched between the simplex with bound $\log Y$ and that with bound $\log Y+\sum_{\ell\in T}\log\ell$. The difference of these volumes is $O_T((1+\log Y)^{d-1})$, proving (*).
\end{proof}

\section{A logarithm saving for every fixed additive fiber}

For
\[
F^+_{a,r}(x)=\#\{q\le x:q\text{ prime},\ A(a,q)=r\},\qquad r\ge3,
\]
if $r\mid a+1$, the fiber has at most one input. Otherwise set
\[
T=\{\ell\le r:\ell\text{ prime},\ \ell\nmid a+1\},\qquad d=|T|.
\]
Then
\begin{equation}
\limsup_{x\to\infty}\frac{F^+_{a,r}(x)}{(\log x)^d}\le C_T. \tag{1}
\end{equation}
Hence
\[
F^+_{a,r}(x)=O_{a,r}((\log x)^{\pi(r)-1}).
\]

\begin{proof}
If $r\mid a+1$ and $r\mid a+q+1$, then $r\mid q$, so primality gives $q=r$. Otherwise take a fiber input $q>r$. Every prime divisor of $a+q+1$ is at most $r$, and none divides $a+1$, since such a divisor would also divide the prime $q$, contradicting $q>r$. Also $r\mid a+q+1$. Therefore
\[
q\mapsto \frac{a+q+1}{r}
\]
injects these inputs into the $T$-smooth integers at most $(x+a+1)/r$. Thus
\[
F^+_{a,r}(x)\le \pi(r)+S_T((x+a+1)/r),
\]
and (*) gives (1). For odd $a$, $2\notin T$; for $a=2$, $3\notin T$. Hence $d\le\pi(r)-1$.
\end{proof}

For example, for $a=r=5$, $T=\{5\}$, so
\[
F^+_{5,5}(x)\le \frac{\log x}{\log5}+O(1).
\]

\section{A \texorpdfstring{$1/\zeta(d)$}{1/zeta(d)} improvement for multiplicative fibers}

Let
\[
F^\times_{a,r}(x)=\#\{q\le x:q\text{ prime},\ M(a,q)=r\}.
\]
The fiber is empty if $a=r$. Otherwise put
\[
T=\{\ell\le r:\ell\text{ prime},\ \ell\ne a\},\qquad d=|T|.
\]
For $d\ge2$,
\begin{equation}
\limsup_{x\to\infty}\frac{F^\times_{a,r}(x)}{(\log x)^d}
\le \frac1{\zeta(d)d!\prod_{\ell\in T}\log\ell}. \tag{2}
\end{equation}

\begin{proof}
The kernel $aq+1$ is $T$-smooth because $a\nmid aq+1$. For $q>a$, suppose its exponent vector has gcd $g>1$, so $aq+1=u^g$. Then $u-1\mid aq$, and $1\le u-1<q$, since $u\le\sqrt{aq+1}<q$. As $a,q$ are distinct primes,
\[
u-1\in\{1,a\}.
\]
Thus every nonprimitive kernel belongs to one of the two families
\[
aq+1=2^g\quad\text{or}\quad aq+1=(a+1)^g,
\]
which contribute only $O_a(\log x)$ possibilities.

For primitive exponent vectors, Mobius inversion and (*) give, with $L=\log Y$,
\[
V_T(L)=\sum_{g\le L/\min_{\ell\in T}\log\ell}\mu(g)(S_T(e^{L/g})-1)
=\frac{C_T}{\zeta(d)}L^d+o(L^d).
\]
The accumulated error is $O(L\log L)$ for $d=2$ and $O(L^{d-1})$ for $d\ge3$. Taking $L=\log(ax+1)$ proves (2). This is an upper-bound coefficient, not an assertion that the prime fiber attains the coefficient.
\end{proof}

\section{The output-5 multiplicative fiber at anchor 2}

For prime $q$,
\[
M(2,q)=5
\]
if and only if either $q=2$, or
\begin{equation}
q=\frac{3^\alpha5^\beta-1}{2},\qquad \alpha,\beta\ge1,\quad \alpha\text{ odd},\quad \gcd(\alpha,\beta)=1. \tag{3}
\end{equation}
Every nonexceptional input satisfies $q\equiv7\pmod{30}$, and
\begin{equation}
\limsup_{x\to\infty}\frac{F^\times_{2,5}(x)}{(\log x)^2}
\le\frac1{3\zeta(2)\log3\log5}. \tag{4}
\end{equation}

\begin{proof}
For odd $q$, the odd 5-smooth kernel is
\[
2q+1=3^\alpha5^\beta,\qquad \beta\ge1.
\]
If $\alpha=0$, then $(5^\beta-1)/2$ is even, so primality forces $q=2$. Hence $\alpha\ge1$. Modulo 4, $\alpha$ is odd. If $g=\gcd(\alpha,\beta)>1$, then $g$ is odd. Putting $t=3^{\alpha/g}5^{\beta/g}\ge15$,
\[
q=\frac{t-1}{2}(1+t+\cdots+t^{g-1}),
\]
contradicting primality. Conversely, (3) directly makes the kernel 5-smooth with greatest prime factor 5. Divisibility by 3 and 5, together with oddness, gives $q\equiv7\pmod{30}$.

For (4), the unrestricted weighted exponent triangle has leading coefficient $1/(2\log3\log5)$. Restricting $\alpha$ to odd values contributes $1/2$, while coprimality over odd common divisors contributes
\[
\sum_{g\text{ odd}}\frac{\mu(g)}{g^2}=\frac4{3\zeta(2)}.
\]
Multiplying gives (4). The formula still requires the resulting $q$ to be prime.
\end{proof}

\section{The output-5 multiplicative fiber at anchor 3}

For prime $q$,
\[
M(3,q)=5
\]
if and only if
\begin{equation}
q=\frac{2^\alpha5^\beta-1}{3},\qquad
\alpha,\beta\ge1\text{ odd},\quad \gcd(\alpha,\beta)=1. \tag{5}
\end{equation}
Every such input satisfies $q\equiv3\pmod{10}$, and
\begin{equation}
\limsup_{x\to\infty}\frac{F^\times_{3,5}(x)}{(\log x)^2}
\le\frac1{6\zeta(2)\log2\log5}. \tag{6}
\end{equation}

\begin{proof}
The input $q=2$ gives output 7. Otherwise the kernel is even, not divisible by 3, and
\[
3q+1=2^\alpha5^\beta,\qquad \alpha,\beta\ge1.
\]
Modulo 3 the exponents have the same parity. If both were even, the kernel would be $t^2$, and
\[
q=\frac{(t-1)(t+1)}3,
\]
with both resulting factors greater than 1, contradicting primality. Hence both exponents are odd. If $g=\gcd(\alpha,\beta)>1$, put $t=2^{\alpha/g}5^{\beta/g}$; then $t\equiv1\pmod3$ and
\[
q=\frac{t-1}{3}(1+t+\cdots+t^{g-1}),
\]
again contradicting primality. The converse is immediate. Oddness and divisibility of $3q+1$ by 5 give $q\equiv3\pmod{10}$.

For (6), the two odd-exponent restrictions contribute $1/4$, and coprimality contributes $4/(3\zeta(2))$.
\end{proof}

\section{Explicit exponential escape and sparse orbits}

Fix a prime anchor $a$ and set $q_{n+1}=E_a(q_n)$. Except for the constant orbit $a=2,q_0=3$,
\begin{equation}
q_n+1\ge2^n(q_0+1). \tag{7}
\end{equation}
Consequently, for $X\ge q_0$,
\begin{equation}
\#\{n\ge0:q_n\le X\}
\le1+\left\lfloor\log_2\frac{X+1}{q_0+1}\right\rfloor. \tag{8}
\end{equation}

\begin{proof}
The prime-exponent theorem gives $E_a(q)\ge2q+1$ for odd prime $q$, except $(a,q)=(2,3)$. For $q=2$, the square-exponent congruence gives $E_a(2)\ge5=2q+1$. A nonexceptional step therefore outputs at least 5 and cannot subsequently enter the exceptional state 3. Hence
\[
q_{n+1}+1\ge2(q_n+1)
\]
at every step. Induction yields (7), and solving (7) for $n$ gives (8).
\end{proof}

The Lean development proves (7), a discrete natural-log form of (8), and the corresponding finite cardinality bound.

\section{Iterated exponential images shrink geometrically}

Let
\[
I_{a,k}=E_a^k(\PP).
\]
Then
\begin{equation}
\#(I_{a,k}\cap[2,X])
\le\pi\!\left(\frac{X+1}{2^k}-1\right)+\ind_{a=2}. \tag{9}
\end{equation}
Thus
\begin{equation}
\limsup_{X\to\infty}\frac{\#(I_{a,k}\cap[2,X])}{\pi(X)}\le2^{-k}. \tag{10}
\end{equation}
Moreover,
\begin{equation}
\bigcap_{k\ge1}I_{a,k}=
\begin{cases}
\{3\},&a=2,\\
\varnothing,&a\ne2.
\end{cases} \tag{11}
\end{equation}

\begin{proof}
If a nonexceptional starting prime has its $k$-th iterate at most $X$, (7) gives
\[
q_0\le\frac{X+1}{2^k}-1.
\]
Counting possible starting primes bounds the number of distinct outputs; the exceptional fixed orbit contributes at most one extra value. This proves (9). The prime number theorem then gives (10).

If a nonexceptional prime $r$ has a $k$-step ancestor, then
\[
r+1\ge2^k(q_0+1)\ge3\cdot2^k,
\]
so its ancestry depth is finite. Only $E_2(3)=3$ has arbitrarily deep ancestry, proving (11).
\end{proof}

The Lean development currently proves (9) and the exact all-depth form of (11). The PNT passage to (10) is not yet formalized in this extension module.

\section{Quantitative rarity of additive-multiplicative collisions}

Fix a prime $a$, put
\[
B_a=a^2+a-1,\qquad R_a=P^+(B_a),
\]
and
\[
T_a=\{\ell\le R_a:\ell\text{ prime},\ \ell\nmid a+1\},\qquad d_a=|T_a|.
\]
Then
\[
D_a(X)=\#\{q\le X:q\text{ prime},\ A(a,q)=M(a,q)\}
\]
satisfies
\begin{equation}
D_a(X)\le C_{T_a}(\log X)^{d_a}+O_a((\log X)^{d_a-1}), \tag{12}
\end{equation}
and therefore
\begin{equation}
\frac{D_a(X)}{\pi(X)}=O_a\!\left(\frac{(\log X)^{d_a+1}}{X}\right). \tag{13}
\end{equation}

\begin{proof}
The polynomial-divisor restriction says every common output $r$ divides $B_a$, so $r\le R_a$ and the additive kernel is $R_a$-smooth. For $q>R_a$, the same exclusion argument as in extension 1 shows that the kernel uses only primes in $T_a$. Thus
\[
D_a(X)\le\pi(R_a)+S_{T_a}(X+a+1).
\]
Apply (*) for (12) and the prime number theorem for (13). The set $T_a$ is nonempty because $B_a\equiv-1\pmod{a+1}$, so its greatest prime divisor cannot divide $a+1$.
\end{proof}

\section{Additive-exponential collisions are confined to a finite region}

If
\[
A(p,q)=E_p(q),
\]
then
\begin{equation}
q\le p, \tag{14}
\end{equation}
except at $(p,q)=(2,3)$. If $p,q$ are odd and $p+q+1$ is composite, then
\begin{equation}
p\ge5q+2. \tag{15}
\end{equation}
For a triple collision,
\begin{equation}
q\le\min\!\left(p,\frac{P^+(p^2+p-1)-1}{2}\right). \tag{16}
\end{equation}

\begin{proof}
Outside the exceptional pair, a common output $r$ satisfies
\[
2q+1\le r\le p+q+1,
\]
which gives (14). If the odd additive kernel $N=p+q+1$ is composite, then its cofactor after division by $r=P^+(N)$ is an odd integer at least 3. Therefore
\[
p+q+1=N\ge3r\ge6q+3,
\]
which is (15). For a triple collision, the additive-multiplicative polynomial restriction gives $r\mid p^2+p-1$, so
\[
2q+1\le r\le P^+(p^2+p-1),
\]
and (16) follows together with (14). The exceptional pair $(2,3)$ is not triple because its multiplicative output is 7 while the additive and exponential outputs are 3.
\end{proof}

If $p\ge3$ and $p^2+p-1$ is prime, no triple collision can occur at anchor $p$: the divisor restriction would force the common output to be $p^2+p-1$, whereas (14) bounds it by $p+q+1\le2p+1<p^2+p-1$.

\section{Complete classification of triple collisions with \texorpdfstring{$q=2$}{q=2}}

For prime $p$,
\begin{equation}
A(p,2)=M(p,2)=E_p(2)\quad\Longleftrightarrow\quad p\in\{2,7\}. \tag{17}
\end{equation}
The common output is 5.

\begin{proof}
The existing second-input-2 collision theorem forces any common output to be 5. Direct calculation gives the two solutions: for $p=2$, the three kernels are all 5; for $p=7$, the kernels are 10, 15, and 50, all with greatest prime factor 5.

Conversely assume $p\ne2$. From $M(2,p)=5$ one obtains $p\equiv1\pmod3$. Since $p$ is odd, $p^2+1$ contains exactly one factor 2 and no factor 3. Thus
\[
p^2+1=2\cdot5^k,\qquad k\ge1.
\]
Modulo 3 forces $k$ even, hence $k\ge2$. Modulo 16 then gives $p^2\equiv1\pmod{16}$.

The additive condition gives
\[
p+3=2^\alpha5^\beta,\qquad \alpha,\beta\ge1,
\]
because 3 does not divide $p+3$. If $\beta\ge2$, then $p\equiv-3\pmod{25}$, so $p^2+1\equiv10\pmod{25}$, contradicting $k\ge2$. Hence $\beta=1$.

Modulo 3, $5\cdot2^\alpha=p+3\equiv1$, so $\alpha$ is odd. Since $p^2\equiv1\pmod{16}$,
\[
p\equiv1,7,9,15\pmod{16}.
\]
In all four cases the 2-adic valuation of $p+3$ is at most 2. Hence $\alpha\le2$, and oddness forces $\alpha=1$. Therefore $p+3=10$, so $p=7$.
\end{proof}

\section{Structural separation from the homogeneous additive GPF family}

Define
\[
H(p,q)=P^+(p+q).
\]

\subsection{Algebraic separation}

There is no magma homomorphism from $H$ into any of $A,M,E$, and no injective magma homomorphism from any of $A,M,E$ into $H$.

\begin{proof}
For every prime $p$,
\[
H(p,p)=P^+(2p)=p,
\]
so every element of $H$ is idempotent. None of $A,M,E$ has an idempotent because $p$ divides none of
\[
2p+1,\qquad p^2+1,\qquad p^p+1.
\]
A homomorphism preserves idempotents, which rules out homomorphisms from $H$ to any of the three new magmas. An injective homomorphism into an idempotent magma would force the source operation to be idempotent, ruling out injective homomorphisms in the reverse direction.
\end{proof}

This algebraic statement is Lean-certified on the branch.

\subsection{Dynamical separation}

For odd prime anchors $a,c$, every orbit of
\[
H_c(q)=P^+(c+q)
\]
is eventually periodic, while every orbit of $E_a$ escapes exponentially. More strongly:
\begin{itemize}
  \item no finite-to-one map $h:\PP\to\PP$ satisfies $h\circ E_a=H_c\circ h$;
  \item no map at all satisfies $h\circ H_c=E_a\circ h$.
\end{itemize}

\begin{proof}
If $q$ is odd, $c+q$ is even and composite, so
\[
H_c(q)\le\frac{c+q}{2}.
\]
For $q=2$, trivially $H_c(2)\le c+2$. Thus every input above $c+2$ strictly decreases, while the finite set of primes at most $c+2$ is forward invariant. Every orbit therefore eventually enters that finite set and is eventually periodic.

An escaping $E_a$-orbit contains infinitely many distinct states. A finite-to-one intertwining map cannot send all of them into one finite $H_c$-orbit. Conversely, $H_c(c)=c$; an intertwining map in the reverse direction would send this fixed point to an $E_a$-fixed point, but no such fixed point exists for odd $a$.
\end{proof}

\section{Remaining directions}

The results above do \emph{not} establish asymptotic equivalences for the actual prime fibers, nor their infinitude, nor a global classification of triple collisions with odd second input. Extensions 1--4 and 7 require a formal weighted-lattice counting layer; the $\zeta$-refinements additionally require Mobius inversion. Equation (10) requires formalizing the prime-counting scaling step from the PNT already available in the project.

\end{document}
